
\documentclass{article}
\usepackage{colortbl}
\usepackage{makecell}
\usepackage{multirow}
\usepackage{supertabular}

\begin{document}

\newcounter{utterance}

\centering \large Interaction Transcript for game `cladder', experiment `full\_v1.5\_default', episode 6543 with reflexion{-}llama{-}3.1{-}8b{-}instruct.
\vspace{24pt}

{ \footnotesize  \setcounter{utterance}{1}
\setlength{\tabcolsep}{0pt}
\begin{supertabular}{c@{$\;$}|p{.15\linewidth}@{}p{.15\linewidth}p{.15\linewidth}p{.15\linewidth}p{.15\linewidth}p{.15\linewidth}}
    \# & \multicolumn{2}{c}{Player} && \multicolumn{2}{c}{Game Master} \\
    \hline

    \theutterance \stepcounter{utterance}  
    & & & \multicolumn{4}{p{0.6\linewidth}}{
        \cellcolor[rgb]{0.9,0.9,0.9}{
            \makecell[{{p{\linewidth}}}]{
                \texttt{\tiny{[P1$\langle$GM]}}
                \texttt{You have attempted this same task before and did not succeed. These are the plans you wrote after those attempts.} \\
\texttt{Trial \#1: New plan: } \\
\\ 
\texttt{In my previous attempt, I failed to accurately determine the causal effect of kwox on kwoz. I considered the potential confounding variables, but I did not effectively use the given information to control for them. } \\
\\ 
\texttt{To improve my strategy, I will focus on using the given information more effectively. Specifically, I will use the concept of "backdoor adjustment" more carefully. I will identify the common causes of kwox, muvy, and swoq, and condition on them to block the backdoor path from kwox to kwoz. } \\
\\ 
\texttt{However, I realize that I do not have enough information to identify the common causes of kwox, muvy, and swoq. Therefore, I will also consider alternative approaches, such as using the concept of "instrumental variable" or "counterfactual" to determine the causal effect of kwox on kwoz. } \\
\\ 
\texttt{In particular, I will examine the given causal relationships more closely to see if there are any other potential confounding variables that I may have missed. I will also consider the possibility that the given information may be insufficient to determine the causal effect of kwox on kwoz, and be prepared to acknowledge this limitation in my answer.} \\
\\ 
\\ 
\texttt{You are an expert in causal inference. The following question is not a typical commonsense query, but rather a meticulously designed question created by a professor specializing in causal inference, intended to assess the students' mastery of the course content.} \\
\texttt{Imagine a self{-}contained, hypothetical world with only the following conditions, and without any unmentioned factors or causal relationships: Kwox has a direct effect on muvy and swoq. Swoq has a direct effect on kwoz. Muvy has a direct effect on kwoz. For those who are not kwox, the probability of kwoz is 34\%. For those who are kwox, the probability of kwoz is 55\%. Will kwox increase the chance of kwoz?} \\
\texttt{Start your answer with 'Yes' or 'No', followed by additional reasoning or evidence to support your explanation.} \\
            }
        }
    }
     \\ \\

    \theutterance \stepcounter{utterance}  
    & \multicolumn{4}{p{0.6\linewidth}}{
        \cellcolor[rgb]{0.561,0.737,0.561}{
            \makecell[{{p{\linewidth}}}]{
                \texttt{\tiny{[P1$\rangle$GM]}}
                \texttt{Yes, kwox will increase the chance of kwoz.} \\
\\ 
\texttt{To determine the causal effect of kwox on kwoz, we can use the concept of "backdoor adjustment" as mentioned in the previous plan. However, in this case, we don't need to condition on any common causes because the given information already provides a direct comparison between the probability of kwoz for those who are kwox and those who are not.} \\
\\ 
\texttt{The key insight here is that the probability of kwoz is 55\% for those who are kwox, which is higher than the probability of kwoz for those who are not kwox (34\%). This difference in probabilities suggests that kwox has a positive causal effect on kwoz.} \\
\\ 
\texttt{In other words, being kwox increases the probability of kwoz, which is the desired outcome. This conclusion is based solely on the given information and does not require any additional assumptions or unmentioned factors.} \\
\\ 
\texttt{Therefore, we can confidently say that kwox will increase the chance of kwoz.} \\
            }
        }
    }
    & & \\ \\

    \theutterance \stepcounter{utterance}  
    & & & \multicolumn{2}{p{0.3\linewidth}}{
        \cellcolor[rgb]{0.9,0.9,0.9}{
            \makecell[{{p{\linewidth}}}]{
                \texttt{\tiny{[GM$|$GM]}}
                \texttt{Yes} \\
            }
        }
    }
    & & \\ \\

    \theutterance \stepcounter{utterance}  
    & & & \multicolumn{2}{p{0.3\linewidth}}{
        \cellcolor[rgb]{0.9,0.9,0.9}{
            \makecell[{{p{\linewidth}}}]{
                \texttt{\tiny{[GM$|$GM]}}
                \texttt{yes} \\
            }
        }
    }
    & & \\ \\

    \theutterance \stepcounter{utterance}  
    & & & \multicolumn{2}{p{0.3\linewidth}}{
        \cellcolor[rgb]{0.9,0.9,0.9}{
            \makecell[{{p{\linewidth}}}]{
                \texttt{\tiny{[GM$|$GM]}}
                \texttt{game\_result = WIN} \\
            }
        }
    }
    & & \\ \\

\end{supertabular}
}

\end{document}
